\section*{Chapitre \ref{ch:mx:why-there-is-a-spread} --- Pourquoi il existe une fourchette}

\begin{solution}{exo:mx:why-there-is-a-spread:1}
$s=2\sqrt{0.0004}=0.04$~: quatre cents.
\end{solution}

\begin{solution}{exo:mx:why-there-is-a-spread:2}
Avec $\delta=\tfrac12$~: $a=99+2\times\tfrac{1.2}{2}=100.2$ et $b=99+2\times\tfrac{0.8}{2}=99.8$~; la fourchette vaut $\mu\times2=0.4$.
\end{solution}

\begin{solution}{exo:mx:why-there-is-a-spread:3}
$\lambda=2/(2\times8)=0.125$, $\beta=8/2=4$, gain $\tfrac12\times2\times8=8$, variance a posteriori $4/2=2$.
\end{solution}

\begin{solution}{exo:mx:why-there-is-a-spread:4}
$\gamma\sigma^2T=0.1$~: $a=50+0.1\,(0.5+3)=50.35$ et $b=50-0.1\,(0.5-3)=50.25$. Étant en position vendeuse, elle veut acheter~: son prix de réserve vaut
$50+0.3=50.3$, et les deux cotations se placent autour de lui, au-dessus de la valeur.
\end{solution}

\begin{solution}{exo:mx:why-there-is-a-spread:5}
La croyance devient l'ancien cours vendeur sur l'échelle unitaire~: $\delta'=0.6$. Nouvelles cotations~: $b=100.00$ et $a=100.38$~; la fourchette se resserre de 0.40 à 0.38 et les deux cotations
remontent.
\end{solution}

\begin{solution}{exo:mx:why-there-is-a-spread:6}
L'initié proportionne ses transactions au bruit~: $\beta=\sigma_u/\sqrt{\Sigma_0}$ double, le rapport signal sur bruit du flux d'ordres est inchangé, et il en va de même de
la part de l'information dans le prix (une moitié). Ce qui change, c'est $\lambda$ (divisé par deux~: un marché plus profond) et le gain de l'initié (doublé).
\end{solution}

\begin{solution}{exo:mx:why-there-is-a-spread:7}
Sur un échantillon (graine 41), la PIN vaut 0.149 avec les fluctuations, 0.139 après normalisation, contre 0.098 sur le même échantillon sans fluctuations (vraie valeur 0.107). La normalisation supprime
le facteur commun, mais aussi le volume supplémentaire des jours informés, que le modèle utilise pour les repérer, et les comptages recalibrés ne suivent plus une loi de Poisson (leur
variance n'est pas leur moyenne), si bien que la vraisemblance est mal spécifiée d'une autre manière.
\end{solution}

\begin{solution}{exo:mx:why-there-is-a-spread:8}
La PIN est une part des transactions issue d'un modèle, non une mesure des pertes informées~: les fluctuations communes de l'activité, et non l'information, peuvent la faire monter (la simulation du chapitre
ajoute 0.15 pour des fluctuations d'écart-type 0.5), et elle ne dit rien du coût que chaque transaction informée fait supporter au teneur de marché. Comparer les coûts de sélection adverse
directement (chapitre~5) avant de conclure.
\end{solution}

\begin{solution}{pb:mx:why-there-is-a-spread:1}
\textbf{1.} Une demi-fourchette constante, des signes de transaction indépendants et une valeur suivant une marche aléatoire~; 0.88 ticks sur l'heure simulée.
\textbf{2.} Les signes de transaction simulés sont positivement autocorrélés (0.29 au décalage~1)~: le rebond est en partie annulé, et l'estimateur reste en deçà
de la fourchette affichée de 1.08 ticks.
\textbf{3.} $a=v+\gamma\sigma^2T(Q/2-q)$, $b=v-\gamma\sigma^2T(Q/2+q)$~; fourchette $\gamma\sigma^2TQ$.
\textbf{4.} Les deux cotations baissent de $\gamma\sigma^2Tq$~: le dealer veut vendre et éviter d'acheter davantage.
\textbf{5.} $\Pr(v_H\mid\text{achat})=\tfrac{(1+\mu)/2}{(1+\mu)/2+(1-\mu)/2}=\tfrac{1+\mu}{2}$, donc $a=v_L+(v_H-v_L)\tfrac{1+\mu}{2}$.
\textbf{6.} 0.3.
\textbf{7.} 148, 19 et 6 transactions.
\textbf{8.} Une moitié et trois quarts.
\textbf{9.} Voir la démonstration de la proposition~: $\beta=1/(2\lambda)$ issu du problème de l'initié, $\lambda=\beta\Sigma_0/(\beta^2\Sigma_0+\sigma_u^2)$ issu du
problème du teneur de marché, d'où $\beta=\sigma_u/\sqrt{\Sigma_0}$, $\lambda=\sqrt{\Sigma_0}/(2\sigma_u)$.
\textbf{10.} La moitié~: $\mathrm{Var}(v\mid p)=\Sigma_0/2$.
\textbf{11.} 0.250, 0.498 et 0.991 (contre 0.25, 0.5 et 1).
\textbf{12.} Les \omterm{def:mx:why-there-is-a-spread:informed}{noise traders}, dont les ordres s'exécutent à des prix déplacés par ceux de l'initié.
\textbf{13.} $\alpha\mu/(\alpha\mu+\varepsilon_b+\varepsilon_s)=24/224=0.107$.
\textbf{14.} 0.006.
\textbf{15.} 0.085, 0.153 et 0.207.
\textbf{16.} \emph{Résultat nommé}~: avec un coût de traitement de 0.05 par transaction et une amplitude de la valeur de 1, la sélection adverse représente $\mu/(\mu+0.1)$ de la fourchette
(une moitié pour $\mu=0.1$, trois quarts pour $\mu=0.3$)~; la PIN estimée les jours sans information passe de 0.006 à 0.085, 0.153 et 0.207 à mesure que l'activité commune
les fluctuations s'amplifient jusqu'à des écarts-types de 0.25, 0.5 et 0.75, et la valeur de 0.107 d'un marché informé est lue comme 0.195 avec le plus grand.
\textbf{17.} Une journée avec plus d'achats et de ventes que d'habitude correspond mieux aux composantes des jours d'événement qu'à la composante sans événement, dont les taux sont fixes~; le modèle
n'a aucun autre moyen d'expliquer la dispersion des totaux quotidiens.
\textbf{18.} Que la composante de la PIN liée à l'illiquidité, et non celle liée à l'information asymétrique, est rémunérée dans la coupe transversale des rendements.
\textbf{19.} Modéliser explicitement le facteur commun (taux de transactions non informées variables dans le temps, chocs symétriques sur le flux d'ordres), et mesurer l'information par ce qui suit
les transactions~: l'impact sur les prix qui persiste (la composante permanente du chapitre~5) plutôt que le volume.
\textbf{20.} Pour une action liquide dont la fourchette vaut un tick, surtout de la sélection adverse et des coûts de traitement nets des rabais~; les \omterm{def:mx:why-there-is-a-spread:inventory}{coûts de détention de l'inventaire} sont faibles parce que
les positions sont renouvelées en quelques secondes.
\end{solution}

\begin{solution}{iq:mx:why-there-is-a-spread:1}
\omterm{def:mx:why-there-is-a-spread:processing}{Coût de traitement des ordres} (indices~: les fourchettes baissent avec les frais et les coûts technologiques~; le rebond de Roll), \omterm{def:mx:why-there-is-a-spread:inventory}{coût de détention de l'inventaire} (les cotations se déplacent à l'encontre des positions du dealer~; les fourchettes s'élargissent
avec la volatilité et la taille des transactions), coût de sélection adverse (les prix vont dans le sens des transactions et y restent~; les fourchettes s'élargissent avant les annonces).
\iqlookfor{les trois noms et une preuve, au niveau du mécanisme, pour chacun.}
\end{solution}

\begin{solution}{iq:mx:why-there-is-a-spread:2}
$\Pr(\text{achat}\mid1)=(1+\mu)/2$, $\Pr(\text{achat}\mid0)=(1-\mu)/2$, donc $\Pr(1\mid\text{achat})=(1+\mu)/2$ et, avec un gain nul, $a=(1+\mu)/2$~; le cours acheteur est
$(1-\mu)/2$ et la fourchette $\mu$. \iqlookfor{la règle de Bayes et la condition de gain nul, énoncées avant le calcul.}
\end{solution}

\begin{solution}{iq:mx:why-there-is-a-spread:3}
$\lambda$ est divisé par deux (l'impact baisse avec la profondeur que procure le bruit), $\beta$ double, le gain de l'initié double, et le prix révèle la même moitié de son
information. \iqlookfor{que le caractère informatif ne dépend pas du bruit, parce que l'initié ajuste son échelle.}
\end{solution}

\begin{solution}{iq:mx:why-there-is-a-spread:4}
Abaisser les deux cotations (skew)~: une plus grande chance de vendre et une moindre chance d'acheter davantage, et éventuellement élargir la fourchette si la position est proche de sa limite~; se couvrir dans un
instrument corrélé s'il en existe un. \iqlookfor{l'idée du prix de réserve et le compromis avec la sélection adverse lorsqu'on applique le skew.}
\end{solution}

\begin{solution}{iq:mx:why-there-is-a-spread:5}
L'autocovariance d'ordre un des variations de prix était positive, et non négative~: les signes de transaction sont autocorrélés (fractionnement des ordres), les prix suivent une tendance sur
l'intervalle d'échantillonnage, ou les données mêlent des cotations de plusieurs plateformes. Utiliser des données de cotation si elles sont disponibles, échantillonner à un pas plus large, ou employer un estimateur qui
n'utilise pas le rebond (Corwin--Schultz, Abdi--Ranaldo~; chapitre~5). \iqlookfor{savoir pourquoi l'estimateur échoue, et pas seulement qu'il échoue.}
\end{solution}

\begin{solution}{iq:mx:why-there-is-a-spread:6}
Décomposer-la~: mesurer la fourchette effective et la fourchette réalisée à un horizon où les variations de prix se sont stabilisées~; leur écart est l'impact sur les prix qui
persiste, la part de sélection adverse~; ou ajuster une décomposition structurelle (Huang--Stoll, MRR) ou un VAR. \iqlookfor{la décomposition effective--réalisée et
le choix de l'horizon, tous deux du chapitre~5.}
\end{solution}
